\section{Summable clearance and exact-mark transport on the physical clock}
\label{sec:ss-transport}
We retain the physical model and the actual event \eqref{eq:v76-event},
including $d\nu_R=d\alpha\,dp/(4\pi)$ and
$d\nu_R^*=c_*^{-1}\one_{Y_R^*}d\nu_R$, where
$c_*=91/(10000\pi)$. Statements involving physical observations below
are only made when both compared configurations are actual regular
records of the same ordered word. The geometric statement comes first
and needs no such physicality of the comparison polygon.

\begin{theorem}[One budget for every center and every flight]
\label{thm:ss-clearance-sum}
For any reference and any $Y\in\mathcal S(Q)$ allowed by
Theorem~\ref{thm:ss-quartic}, define
\[
 \kappa_i=\sup_{z\in\Lambda}
 \left|\operatorname{dist}(z,[Y_i,Y_{i+1}])-
              \operatorname{dist}(z,[Q_i,Q_{i+1}])\right|.
\]
Then
\begin{equation}\label{eq:ss-clearance-sum}
 \sum_{i=0}^{m-1}\kappa_i^4\le2CE,\qquad
 \#\{i:\kappa_i\ge\rho\}\le2CE\rho^{-4}\quad(\rho>0).
\end{equation}
The same estimates apply to the absolute differences of the
unvisited-center clearances $\Gamma_i$ in
Section~\ref{sec:sc-clearance}. All centers, including a center
visited on a different flight, and all possible feet on the closed
segments are included. Negative clearance of a nonphysical $Y$ is
permitted and is not suppressed.
\end{theorem}
\begin{proof}
Using the same affine parameter on the two segments gives their
Hausdorff distance at most $\max\{|\delta_i|,|\delta_{i+1}|\}$.
Distance to any fixed point is $1$-Lipschitz for Hausdorff distance,
so the bound holds simultaneously for every lattice center, without
any enumeration. Thus
$\kappa_i^4\le|\delta_i|^4+|\delta_{i+1}|^4$.
The full sum is at most twice the contact sum in
\eqref{eq:ss-sublevel-sum}; each endpoint is in fact counted only
once. Taking infima over the identical set
$\Lambda\setminus\{O_i,O_{i+1}\}$ preserves this uniform difference
bound for $\Gamma_i$. Finally
$\rho^4\#\{i:\kappa_i\ge\rho\}\le\sum_i\kappa_i^4$.
\end{proof}

\subsection{Physical reference class and the original section}
Fix $R$ and $m$. Let $\mathcal C_{R,m}$ be the union of actual regular
word cells whose disk-product minimum is an actual regular
normal-endpoint record $Y^w$. Such a minimum is unique by the retained
critical result, or Corollary~\ref{cor:ss-critical-closure}.
The sufficient strict physicality certificate of
Theorem~\ref{thm:sc-physical} yields a subclass; no claim is made that
all words lie in $\mathcal C_{R,m}$. Word cells and their original
exact-event subsets are Borel, with the usual countable
singularity partition. For this fixed $R,m$, the assignment of the
constant record $Y^w$ on each word cell is therefore Borel; it involves
no unproved measurable choice of auxiliary minima.

For $Q\in\mathcal C_{R,m}$ compare its $m+1$ collision states
$x_i=(\alpha_i,p_i)$ with $\widehat x_i$ of $Y^w$. Angles are taken
in $\mathbb R/(2\pi\mathbb Z)$, and
\[
 d_\infty(x,\widehat x)=
 \max\{\operatorname{dist}_{\mathbb S^1}(\alpha,\widehat\alpha),
                                      |p-\widehat p|\}.
\]
At the last collision, the tangential momentum can be computed from
the incoming flight: specular reflection leaves that component
unchanged. Consequently the terminal state, although its next flight
is not in the record, is fully determined for the comparison.

\begin{theorem}[Full phase budget and section discrepancy]
\label{thm:ss-section}
Put $R_-=9/20$ and
\begin{equation}\label{eq:ss-phase-constant}
 K=128/\tau_-+24C/R_-^4.
\end{equation}
For every actual $Q\in\mathcal C_{R,m}$, writing $E=E_R(Q)$,
\begin{equation}\label{eq:ss-phase-sum}
 \sum_{i=0}^m d_\infty(x_i,\widehat x_i)^4\le KE.
\end{equation}
Let $\eta_R=\one_{Y_R^*}$ denote the original section indicator, not a
smoothed or enlarged section, and define
\begin{align*}
 N_{\rm cut}(Q)&=\sum_{i=0}^m|\eta_R(x_i)-\eta_R(\widehat x_i)|,\\
 B_\rho(Q)&=\sum_{i=0}^m
  \one_{\{\operatorname{dist}_\infty(x_i,\partial Y_R^*)\le\rho\}}.
\end{align*}
Then for every $\rho>0$
\begin{equation}\label{eq:ss-section-bound}
 N_{\rm cut}(Q)\le B_\rho(Q)+K\rho^{-4}E_R(Q).
\end{equation}
Both endpoints are included once, whereas the original occupation
count $A_m=\sum_{i=0}^{m-1}\eta_R(x_i)$ still excludes the final
state. In particular
$|A_m(Q)-A_m(Y^w)|\le N_{\rm cut}(Q)$.
\end{theorem}
\begin{proof}
Since $Y^w$ is physical and lies in $\mathcal S(Q)$, all its vertices
are on the original disks. The unit normals differ by
$|\widehat n_i-n_i|=|\delta_i|/R$. For angles at circular distance
$\theta\in[0,\pi]$, $2\sin(\theta/2)\ge2\theta/\pi$.
Thus
\[
 \operatorname{dist}_{\mathbb S^1}(\alpha_i,\widehat\alpha_i)
 \le\frac{\pi}{2R}|\delta_i|\le\frac2{R_-}|\delta_i|.
\]
Tangential projection of the outgoing unit vector gives
$|p_i-\widehat p_i|\le|e_i|+|\delta_i|/R_-$ for $i<m$.
The terminal incoming projection gives the same estimate with
$e_{m-1}$. Using $(u+v)^4\le8(u^4+v^4)$ and bounding a maximum
by a sum, we get
\[
 d_\infty(x_i,\widehat x_i)^4
 \le8|e_{\iota(i)}|^4+24R_-^{-4}|\delta_i|^4,
 \quad\iota(i)=\min(i,m-1).
\]
The last flight index is used twice and the others once. Since the
difference of two unit vectors has norm at most $2$,
\[
 \sum_i|e_i|^4\le4\sum_i|e_i|^2
 \le 8G/\tau_-\le8E/\tau_-.
\]
Summing the phase estimate and using
\eqref{eq:ss-sublevel-sum} proves \eqref{eq:ss-phase-sum}.

The original section is a finite union of open rectangles in these
circle--momentum coordinates. If its indicator differs at two points,
a shortest circle--momentum path between them meets its boundary.
Hence a discrepancy at index $i$ implies
$\operatorname{dist}_\infty(x_i,\partial Y_R^*)
 \le d_\infty(x_i,\widehat x_i)$.
If the left side is at most $\rho$, it is counted by $B_\rho$;
otherwise its indicator is at most
$\rho^{-4}d_\infty(x_i,\widehat x_i)^4$.
Sum over all $i$ and use \eqref{eq:ss-phase-sum}. Points on the
boundary are explicitly counted, not assigned zero cost by a bulk-null
argument. The occupation assertion follows by dropping the last
nonnegative discrepancy.
\end{proof}

\begin{proposition}[Exact lattice labels and all four twist variables]
\label{prop:ss-twists}
On $\mathcal C_{R,m}$ the true count $m$ and lattice displacement $k$
of $Q$ and $Y^w$ agree. Their true roofs satisfy
$0\le L(Q)-L(Y^w)\le E_R(Q)$. For every
$\xi\in\mathbb R^2$, $s,u\in\mathbb R$,
\begin{align}\label{eq:ss-twist}
 &\left|e^{i[\xi\cdot k+sL(Q)+uA_m(Q)]}
         -e^{i[\xi\cdot k+sL(Y^w)+uA_m(Y^w)]}\right|\notag\\
 &\hspace{15mm}\le |s|E_R(Q)+2|\sin(u/2)|N_{\rm cut}(Q).
\end{align}
The factors $\eta_R(x_0)\eta_R(x_m)$, if included, incur in addition
at most the two endpoint discrepancies. No label stability is asserted
when that discrepancy has not been paid.
\end{proposition}
\begin{proof}
The physical ordered centers are identical; thus the displacement and
number of flights are identical, without dropping a reflection.
The actual critical minimum has length $\ell_w(R)$, so the inherited
endpoint-gap theorem gives the roof inequality. The triangle
inequality and $|e^{ia}-e^{ib}|\le|a-b|$ bound the roof part.
For an integer $q$, telescoping $q$ unit phase increments gives
$|e^{iuq}-1|\le |q|\,|e^{iu}-1|=2|q||\sin(u/2)|$,
including negative $q$. Take $q=A_m(Q)-A_m(Y^w)$ and use
Theorem~\ref{thm:ss-section}. A product of two indicators has difference
at most the sum of the differences of its two factors. This also
shows exactly where the initial and terminal section tests enter.
\end{proof}

\begin{theorem}[Positive discrepancy payment on the unchanged roof]
\label{thm:ss-source}
For every $R,m,n,k$ and nonnegative Borel weight $h$, let $\Pi$ be the
original exact event \eqref{eq:v76-event} and define, on the class
$\mathcal C_{R,m}$ only,
\[
 \mathcal M_h[f]=(L_R)_*
       (\one_{\mathcal C_{R,m}}\one_\Pi h f\,d\nu_R^*).
\]
The pushforward is always by the actual original roof $L_R(Q)$,
not the comparison roof. As an inequality of positive measures,
\begin{equation}\label{eq:ss-source-payment}
 \mathcal M_h[N_{\rm cut}]
 \le \mathcal M_h[B_\rho]+K\rho^{-4}\mathcal M_h[E_R]
 \quad(\rho>0).
\end{equation}
It remains valid after a countable sum over all physical word cells
and nonnegative marked copies, with their actual multiplicities. If
all three measures have densities, the same inequality holds between
the densities almost everywhere, at the same $R,m,n,k,t$.
\end{theorem}
\begin{proof}
Multiply \eqref{eq:ss-section-bound} by the nonnegative actual weight
$\one_{\mathcal C_{R,m}}\one_\Pi h$ and integrate its product with
any nonnegative Borel test of $L_R(Q)$. This proves the measure
inequality, first for bounded weights and then by monotone
convergence. Countable marked sums are handled by Tonelli before any
supremum or limiting operation. The conditional collision probability
already contains exactly the factor $c_*^{-1}$; no second section
normalization is introduced. There is no coordinate change in this
proof, and hence no unrecorded Jacobian or physical time multiplier.
The density statement is the order property of Radon--Nikodym
derivatives when absolute continuity has actually been established.
\end{proof}

\begin{remark}[The exact remaining normalization]\label{rem:ss-open}
The new term $B_\rho$ is a specified observation-boundary occupation,
and $E_R$ is the actual endpoint defect. Their occurrence is an
explicit error accounting for real labels, not a claim that either
has small normalized density. The class complementary to
$\mathcal C_{R,m}$ is not discarded. Equations
\eqref{eq:ss-source-payment} and \eqref{eq:ss-twist} alone give no
$m^{-2}$ estimate, no high-$s$ uniform decay, and no comparison-pool
capacity. In particular choosing $\rho$ as a function of $m$ requires
estimates uniform at those widths; it is not justified by a
fixed-width result. The singular paired endpoint traces are not bulk
pushforwards and are not declared paid by this theorem. The complete
source must still satisfy the original
\eqref{eq:v112-remaining}, with all exact labels, roof intervals,
ordered suprema and outermost $B$ limit unchanged. This paper does
not assert completion of the raw pointwise theorem or its
four-coordinate inversion consequence.
\end{remark}
